\section{Governance and methodology evolution}
\label{sec:governance}

\subsection{Versioning}

The methodology version --- \methver{} --- stamps this document, the
Methodology Report in the repository, and every configuration file, and is to
be recorded in every assessment result the engine produces. The version is a
date stamp rather than a semantic version, because a methodology does not have
a meaningful notion of a backward-compatible change: any alteration to a value
may alter a ranking.

A change to any methodology value requires a version bump and a corresponding
update to this document in the same change set. Continuous integration enforces
the mechanical half of this --- version lock-step across configuration files,
and citations for every performance value (\Cref{sec:evidence-rules}) --- while
the substantive half, that the documentation actually describes the new value,
is enforced in review.

Existing assessments are never to be silently recomputed: a result retains the
version that produced it, the interface indicates when a newer methodology
version is available, leaving the decision to re-run with the user, and
comparisons between results produced under differing versions raise a warning.

Version \methver{} is the first release to exercise this properly, because it
is the first to break something. The intervention field became a list and the
input schema gained four fields, so stored \emph{drafts} cannot be read
unchanged. They are migrated \textbf{explicitly} rather than silently: the
single intervention is wrapped into a one-element list, the four new fields are
left absent rather than defaulted, and the migration is itemised for the user.
The migration notes distinguish two cases that would otherwise be confused ---
a slug that no longer names any selectable entry, whose draft must be
re-chosen, and a slug that still names a catalogue entry but whose inherited
evidence class changed, which needs no re-selection. Telling a user to
re-select an intervention that is still there would be a false alarm about
their own saved work. Stored \emph{results} are not touched in either case.

\subsection{The public change process}

Methodology changes are proposed as public pull requests, accompanied by the
evidence supporting them. A proposal that changes a quantitative value is
expected to state which source supports the new value, what was verified about
that source and how, and why the existing value should be displaced. The
verification standard is the one described in \Cref{app:verification}:
bibliographic details checked against the publisher record, an institutional
repository, or an indexing service, with the verification level recorded
honestly rather than overstated.

\begin{keynote}
\noindent\textbf{A design may be approved; a citation is only settled once it
has been read.} Version \methver{} supplies the case that makes this concrete.
An envelope floor was approved on the strength of a figure attributed to
\textcite{yao2023a}; verification before the value entered configuration found
that the paper reports no such figure and in fact measures an order of
magnitude less by day, and the floor moved (\Cref{sec:water-body}). Four
further citation corrections were made in the same pass without moving a value.

The process rule that follows is stated here rather than left as folklore:
\textbf{approval of a design is not approval of the citations inside it}, and
verification runs after approval and before configuration, with the authority
to change the approved value.
\end{keynote}

Because the methodology lives in configuration rather than in code, the diff of
such a proposal is legible to a domain expert who does not read the
implementation language. This is a deliberate property. A reviewer evaluating a
change to the green fa\c{c}ade envelope should be able to see a temperature
range, a confidence rating, and a citation change, and nothing else.

\subsection{How to challenge this methodology}

Methodology critique is a first-class contribution to this project, not an
inconvenience to be managed. The route is to open an issue in the repository
at \url{\repourl}, referencing the section number of this document and, where
possible, the literature supporting the alternative position.

Six areas are in particular need of external scrutiny, and are listed in
descending order of how much a change would affect the tool's outputs.

\begin{enumerate}
  \item \textbf{The aggregation weights} (\Cref{sec:aggregation}), and
  specifically the equity-forward double contribution of vulnerability
  (\Cref{sec:effective-weights}). This is the tool's most consequential
  unsourced choice and a value position rather than a finding. The
  sensitivity analysis of \Cref{sec:sensitivity-results} quantifies its
  influence: rank orderings prove substantially stable under single-weight
  perturbation, which bounds --- but does not settle --- the dispute.

  \item \textbf{The green fa\c{c}ade and bioretention envelopes}
  (\Cref{sec:weak-points}). For green fa\c{c}ades the sources conflict outright
  and the tool adopts a conservative resolution based on a hypothesis about
  measurement position that neither source states; eleven catalogue entries
  inherit the result. For rain gardens and bioswales no source quantifies
  air-temperature cooling at all, and the adopted range is reasoned from
  adjacent evidence. Both are candidates for replacement by better evidence,
  and both would benefit from a reviewer who knows this literature better than
  the derivation in \Cref{app:typologies} demonstrates.

  \item \textbf{The two derived productive archetypes, and the width of the
  large water body envelope} (\Cref{sec:derived-archetypes,sec:water-body}).
  The productive classes are bounded from adjacent evidence rather than
  measured, and the tool can cite the absence \parencite{kumar2024} but cannot
  fill it. The water envelope spans a factor of thirty because remote sensing
  and field measurement disagree by an order of magnitude; a reviewer who can
  resolve that disagreement on air temperature would narrow the widest range in
  the library.

  \item \textbf{The non-adoption of \textcite{kumar2024}'s higher figures}
  (\Cref{sec:non-adoption,sec:limit-non-adoption}). Adopting them selectively
  would break internal consistency and adopting them wholesale is a full
  recalibration with its own sensitivity analysis. This is a live disagreement
  with a recent review, held open deliberately, and an argument for closing it
  either way is welcome.

  \item \textbf{The transferability of the energy sensitivity}
  (\Cref{sec:energy}). The \SIrange{2}{4}{\percent} per \si{\degreeCelsius}
  figure is applied globally from predominantly North American evidence.
  Reviewers with access to equivalent evidence for other building stocks and
  climates would materially improve this part of the methodology, and evidence
  that the figure does \emph{not} transfer would be equally valuable.

  \item \textbf{The downward correction for modelling bias}
  (\Cref{sec:modelling-bias}). The correction rests on a single observable
  ratio for a single measure, used to justify a direction rather than a
  magnitude. A reviewer who considers the resulting envelopes too conservative
  --- or not conservative enough --- is disputing exactly the right thing.
\end{enumerate}

\subsection{What a good challenge looks like}

The most useful critiques share a shape. They identify a specific value or
rule, by section number; they state what is wrong with it, distinguishing
between an arithmetic error, a misread source, a source that does not support
the weight placed on it, and a value judgment the critic rejects; and where the
critique is evidential rather than a difference of values, they name a source
and what it actually reports, including which metric it reports it in.

The last point is the one most often missed. A critique of a cooling envelope
that cites a figure without stating whether it is air temperature, land surface
temperature, or a comfort index cannot be acted on, because the methodology
cannot adopt a value whose metric is unknown (\Cref{sec:metric}). This is the
same standard the methodology holds itself to.

Critiques that reject a value judgment --- the equity weighting being the
obvious case --- are welcome and will be recorded, but they will generally be
resolved by making the configuration easier to rebalance and the consequences
easier to see, rather than by replacing one contested judgment with another.
